
\subsection{Les opérateurs de base}

\begin{frame}[fragile]{Les opérateurs de base}

Nous allons voir tous les opérateurs permettant d'exécuter les
requêtes mono-table.

\begin{minted}[frame=leftline,
               baselinestretch=.9,
               fontsize=\footnotesize,
               framesep=2mm]{sql}
select a1, a2, ..., an 
from T 
where condition
\end{minted}

\vfill

Pour accéder à une table, deux solutions (pas plus!)

\begin{itemize}
 \item \alert{Accès séquentiel}. Tous les tuples de la table 
      sont examinés, dans l'ordre du stockage.
 \item \alert{Accès par adresse}. Si on connaît l'adresse du
      tuple, on peut aller lire
      directement le bloc et obtenir ainsi un accès optimal.
\end{itemize}

\vfill


\end{frame}


\begin{frame}[fragile]{Parcours de table et d'index}

\alert{Parcours séquentiel}: déjà vu avec l'opérateur \texttt{FullScan}.
\smallskip

\alert{Efficacité}: pas si mal, si en séquence, sur une table compacte.

\vfill

\alert{Parcours d'index} (index = arbre B). Deux phases pour l'opérateur \texttt{IndexScan}:

\begin{itemize} 
  \item \alert{Parcours de la racine vers la feuille} (pendant le \texttt{open()})
  \item \alert{Parcours en séquence des feuilles} (une étape par \texttt{next()})
  
\end{itemize}

\smallskip

\alert{Efficacité}: très efficace, quelques lectures logiques (index en mémoire)

\vfill

Lequel préférer (quand on a le choix)? En général, l'index, mais pas toujours.

\end{frame}

%%%%
\begin{frame}[fragile]{Accès par adresse}

Opérateur \texttt{DirectAccess}: s'appuie sur une source (un parcours d'index).

\vfill

\begin{minted}[frame=leftline,
               baselinestretch=.9,
               fontsize=\footnotesize,
               framesep=2mm]{bash}
   function nextDirectAccess
    {
      # Prenons l'adresse de l'enregistrement a lire
      $a = $source.next();
      
      # Plus d'adresse? On renvoie null
      if ($a = null) then
        return null;
      else
        # On effectue par une lecture (logique) du bloc
        $b = $GA.lecture ($a.adresseBloc);      
        # On prend l'enregistrement dans le bloc
        $e = $b.get ($a.adresseLocale)           
        return $e;
      fi 
     }
\end{minted}

\end{frame}


%%%%
\begin{frame}[fragile]{Opérateur de sélection}

Très simple: filtre les tuples fournis par la source.

\begin{minted}[frame=leftline,
               baselinestretch=.9,
               fontsize=\footnotesize,
               framesep=2mm]{bash}
    function nextFilter
    {
       # On prend un tuple de la source
      $tuple = $source.next();
      # On continue tant que la condition n'est pas satisfaite, 
      # ou la source parcourue
      while ($tuple != null  and $tuple.test($C) = false)
      do
        $tuple = $source.next();        
      done
      
      return $tuple;
    }
\end{minted}

\end{frame}

\begin{frame}[fragile]{Plan d'exécution pour requête mono table}

\begin{minted}[frame=leftline,
               baselinestretch=.9,
               fontsize=\footnotesize,
               framesep=2mm]{sql}
select titre from Film where genre='SF' and annee = 2014
\end{minted}

\medskip

Les solutions:

\figSlideWithSize{planEx-monotable}{10cm}

\end{frame}


\begin{frame}[fragile]{Index ou pas index?}

Un cas extrême!

\medskip

\figSlideWithSize{interBtree}{10cm}

\medskip

En pratique: le SGBD décide en fonction des statistiques.

\end{frame}


\begin{frame}[fragile]{Pour compléter}

Un plan à exécuter avec index.

\begin{minted}[frame=leftline,
               baselinestretch=.9,
               fontsize=\footnotesize,
               framesep=2mm]{sql}
    select * from   Film
    where  idFilm = 20 
    and    titre = 'Vertigo'
\end{minted}

\vfill

Un plan à exécuter sans index.

\begin{minted}[frame=leftline,
               baselinestretch=.9,
               fontsize=\footnotesize,
               framesep=2mm]{sql}
    select * from   Film
    where  idFilm = 20 
    or    titre = 'Vertigo'
\end{minted}

\end{frame}